Идея bagging (bootstrap aggregation) заключается в следующем.
Пусть обучающая выборка состояла из $n$ объектов.
Выберем из неё $n$ примеров равновероятно, с возвращением.
Получим новую выборку $X_1$, в которой некоторых элементов исходной выборки не будет, а какие-то могут войти
несколько раз.
С помощью некоторого алгоритма $b$ обучим на этой выборке модель $b_1(x) = b(x, X_1)$.
Повторим процедуру: сформируем вторую выборку $X_2$ из $n$ элементов с возвращением и с помощью того же алгоритма
обучим на ней модель $b_2(x) = b(x, X_2)$.
Повторив процедуру $k$ раз, получим $k$ моделей, обученных на $k$ выборках.
Чтобы получить одно предсказание, усредним предсказания всех моделей:
\[ a(x) = \frac{1}{k}\bigl(b_1(x) + \dots + b_k(x)\bigr) \]

\subsection{Мотивация через bias–variance}

Если предположить, что базовые алгоритмы некоррелированы, то:

\begin{align*}
    V_X[a(x,X)] &= \frac{1}{k^2}\sum_{i=1}^k V_X b(x,X_i) \\
    &= \frac{1}{k^2}\sum_{i=1}^k V_X b(x,X) \\
    &= \frac{1}{k}\, V_X b(x,X).
\end{align*}

Получилось, что в этом случае дисперсия композиции в $k$ раз меньше дисперсии отдельного алгоритма.

\subsection{Bootstrap}

Процесс генерации подвыборок с помощью семплирования с возвращением называется (bootstrap), а модели $b_1(x),
\dots,b_k(x)$ часто называют базовыми алгоритмами (хотя, наверное, лучше было бы назвать их базовыми моделями).
Модель $a(x)$ называется ансамблем этих моделей.

Посмотрим, что происходит с качеством предсказания при переходе от одной модели к ансамблю.
Сначала убедимся, что смещение ансамбля не изменилось по сравнению со средним смещением отдельных моделей.
Будем считать, что когда мы берём матожидание по всем обучающим выборкам $X$, то в эти выборки включены также все
подвыборки, полученные by bootstrap.

\begin{align*}
    \operatorname{bias}_X a(x,X) &= f(x) - \mathbb{E}_X[a(x,X)] \\
    &= f(x) - \mathbb{E}_X\!\Bigl[\frac{1}{k}\sum_{i=1}^k b(x,X_i)\Bigr] \\
    &= f(x) - \frac{1}{k}\sum_{i=1}^k \mathbb{E}_X[b(x,X_i)] \\
    &= f(x) - \frac{1}{k}\sum_{i=1}^k \mathbb{E}_X[b(x,X)] \\
    &= f(x) - \mathbb{E}_X b(x,X) \\
    &= \operatorname{bias}_X b(x,X).
\end{align*}

Получили, что смещение композиции равно смещению одного алгоритма.
Теперь посмотрим, что происходит с разбросом.

\begin{align*}
    V_X[a(x,X)] &= \mathbb{E}_X\bigl[a(x,X) - \mathbb{E}_X[a(x,X)]\bigr]^2 \\
    &= \mathbb{E}_X\!\Bigl[\frac{1}{k}\sum_{i=1}^k b(x,X_i) - \mathbb{E}_X\!\Bigl[\frac{1}{k}\sum_{i=1}^k b(x,X_i)
    \Bigr]\Bigr]^2 \\
    &= \frac{1}{k^2}\,\mathbb{E}_X\!\Bigl[\sum_{i=1}^k \bigl(b(x,X_i) - \mathbb{E}_X b(x,X_i)\bigr)\Bigr]^2 \\
    &= \frac{1}{k^2}\sum_{i=1}^k \mathbb{E}_X\bigl(b(x,X_i) - \mathbb{E}_X b(x,X_i)\bigr)^2 \\
    &\quad + \frac{1}{k^2}\sum_{k_1\neq k_2} \mathbb{E}_X\!\Bigl[\bigl(b(x,X_{k_1}) - \mathbb{E}_X b(x,X_{k_1})
    \bigr)\bigl(b(x,X_{k_2}) - \mathbb{E}_X b(x,X_{k_2})\bigr)\Bigr] \\
    &= \frac{1}{k^2}\sum_{i=1}^k V_X b(x,X_i) + \frac{1}{k^2}\sum_{k_1\neq k_2} \operatorname{cov}
    \bigl(b(x,X_{k_1}), b(x,X_{k_2})\bigr).
\end{align*}

\subsection{Метод случайных подпространств}

Суть в том, что в каждой вершине случайно выбираются $n<N$ признаков, где $N$ -- полное число признаков
(метод случайных подпространств), и среди них ищется оптимальный сплит.
Такой приём как раз позволяет управлять степенью скоррелированности базовых алгоритмов.

\subsection{Случайный лес}

В предыдущем разделе мы сделали предположение, что базовые алгоритмы некоррелированы, и за счёт этого получили очень
сильное уменьшение дисперсии у ансамбля относительно входящих в него базовых алгоритмов.
Однако в реальной жизни добиться этого сложно: ведь базовые алгоритмы учили одну и ту же зависимость на
пересекающихся выборках.
Поэтому будет странно, если корреляция на самом деле нулевая.
Но на практике оказывается, что строгое выполнение этого предположения не обязательно.
Достаточно, чтобы алгоритмы были в некоторой степени непохожи друг на друга.
На этом строится развитие идеи bagging для решающих деревьев -- случайный лес.

Построим ансамбль алгоритмов, где базовый алгоритм -- это решающее дерево.
Будем строить по следующей схеме:

\begin{enumerate}
    \item Для построения $i$-го дерева

    \begin{itemize}
        \item Сначала, как в обычном bagging, из обучающей выборки $X$ выбирается с возвращением случайная подвыборка
        $X_i$ того же размера, что и $X$.
        \item В процессе обучения каждого дерева используется метод случайных подпространств
    \end{itemize}

    \item Чтобы получить предсказание ансамбля на тестовом объекте, усредняем отдельные ответы деревьев
    (для регрессии) или берём самый популярный класс (для классификации).
    \item Мы построили Random Forest (случайный лес) -- комбинацию bagging и метода случайных подпространств над
    решающими деревьями.
\end{enumerate}
При построении случайного леса есть несколько степеней свободы